% ==============================================================
%  Chapter 11: 50 Essential Q&As for Developers and Professionals
% ==============================================================
\chapter{50 Essential Q\&As for Developers and Professionals}
\label{ch:dev-qa}

These questions represent the most common, most important, and most misunderstood
topics when professionals start integrating AI into their technical work.
The answers are deliberately direct and practical.

\qa{How should I structure prompts for coding tasks?}
{Always provide: (1) language and version, (2) framework if relevant, (3) what
the function receives as input, (4) what it should return, (5) constraints such as no
external libraries or style guide compliance.
One clear goal per prompt is more effective than a multi-part request.
State the output format you want --- code only, code with explanation, or code with
tests.}

\qa{Can I trust AI-generated code in production?}
{Treat it the way you would treat code from a capable junior developer: review it,
test it, check edge cases, and confirm it handles your specific business rules.
AI is right most of the time on common patterns and wrong more often on
domain-specific or context-sensitive logic.
Never push AI-generated code to production without review and test coverage.}

\qa{How do I debug effectively with AI?}
{Don't just paste the error message.
Provide the full stack trace, the relevant code, what you expected to happen, what
actually happened, and what you have already tried.
Ask for the three most likely root causes and how to diagnose each --- not just a fix.
Reasoning diagnostically gives better results than asking for a direct patch.}

\qa{How do I give AI enough context about my codebase?}
{Paste only the relevant files or functions --- not the entire repository.
Summarise your architecture in 2--3 sentences at the top of the prompt.
For large codebases, use tools like Cursor or GitHub Copilot Workspace, which index
your repository and retrieve relevant context automatically via embeddings.}

\qa{Should I use AI for security-critical code?}
{Use AI to understand patterns and generate boilerplate, but manually review all
authentication, authorisation, encryption, and input validation code.
AI regularly produces plausible-looking security code with subtle flaws.
Run a static security analyser (Semgrep, Bandit, ESLint security plugins) on any
AI-generated code regardless, and treat the result as a starting point.}

\qa{How do I avoid AI hallucinations in code?}
{Test the code rather than reading it.
Verify all library method names and signatures against official documentation.
Ask AI ``how would I verify this is correct?'' to surface its own uncertainty.
For libraries released or updated after the model's training cutoff, hallucinations
in API details are especially common.}

\qa{How do I use AI to understand legacy code?}
{Paste the function and ask: ``Explain this code line by line, then summarise what
business problem it solves and what would break if it were removed.''
AI is excellent at reverse-engineering intent from code.
Ask for a dependency summary or call graph to understand broader coupling.}

\qa{Can AI write unit tests I can rely on?}
{AI generates test structure and happy-path cases quickly, but misses edge cases
unless you ask for them.
First, ask AI to list all cases to test.
Review that list, add anything missing, and then ask for the test code.
Review tests as carefully as you review production code --- a test that doesn't
catch a real regression is worse than no test.}

\qa{What is the difference between a chat AI and a coding assistant?}
{Chat AI (ChatGPT, Claude) is better for reasoning, explanation, architecture
discussion, and learning --- you work conversationally.
Coding assistants (Copilot, Cursor) are better for inline completion while you type
and for keeping you in flow.
Use both: chat AI for the thinking, coding assistants for the implementation.
They are complementary, not alternatives.}

\qa{How do I maintain my own skills while using AI?}
{Deliberately solve problems without AI at regular intervals.
Before accepting AI output, try to predict what the answer should be.
Ensure you can explain and modify every line of AI-generated code you commit.
Use AI to accelerate work you already understand --- not to skip understanding
entirely.}

\qa{How do I use AI for refactoring?}
{Paste the code and say: ``Refactor this for readability. Apply [principle/pattern].
Explain every change you make.''
Then compare the before and after carefully.
Ask ``does this refactor change any behaviour?'' to catch semantic drift.
AI refactoring is reliable for obvious patterns but can introduce subtle bugs in
complex conditional logic.}

\qa{Can AI help with architecture decisions?}
{AI is excellent for generating options and tradeoffs.
Provide your constraints explicitly: team size, skill set, existing services, budget,
latency requirements, and compliance needs.
Ask for three approaches with tradeoffs rather than one recommendation.
AI cannot know your organisation's risk appetite or team dynamics --- you supply those.}

\qa{How do I use AI for performance optimisation?}
{Profile first to confirm the actual bottleneck before asking AI to optimise anything.
Paste the slow function and ask: ``Identify performance bottlenecks and suggest
optimisations with their big-O complexity impact.''
Always benchmark before and after to verify the improvement is real.}

\qa{How do I use AI for regex?}
{Describe what you want to match in plain English and give 3--4 concrete examples
of input strings that should match and strings that should not.
Ask AI to explain the regex pattern it generates alongside the code.
Always test against your actual data --- edge cases and encoding issues are common.}

\qa{How do I use AI for SQL queries?}
{Give the table names, column names, and data types.
Describe the result you want.
Ask for an explanation of the JOIN strategy and whether the query is index-friendly.
Review any user-supplied input handling carefully for SQL injection risk.}

\qa{How do I use AI for shell scripting?}
{Describe the task step by step.
Ask for explanations of each command.
Ask: ``What could go wrong with this script, and what safeguards should I add?''
Test in a non-production environment first.
AI shell scripts often lack error handling --- ask for it explicitly.}

\qa{How do I verify AI suggestions before committing?}
{Run the code and its tests.
Read the diff carefully --- every line, not just the changed lines.
Search for any new library calls in official documentation to verify accuracy.
Run a linter and static analyser as a final pass.
If you cannot explain a line, that is a signal to investigate before committing.}

\qa{Can AI help with concurrency and threading problems?}
{AI knows common patterns (mutex, semaphore, actor model, CSP) and can explain
them clearly.
It misses race conditions in complex real-world scenarios regularly.
Treat any multithreaded code AI generates as a starting point requiring careful review.
Ask specifically: ``What are the race conditions possible in this code?''}

\qa{What is the best way to use AI in a team setting?}
{Agree on what categories of code require human review (all production code, at
minimum).
Share effective prompt templates in your team wiki.
Don't use AI to bypass your code review process --- use it to make code review
faster by providing better-quality first drafts.}

\qa{How do I use AI for onboarding to a new codebase?}
{Paste key files and ask: ``What does this service do, what are its main components,
and what is the entry point?''
Follow up: ``If I needed to add a new endpoint, what files would I need to change
and in what order?''
AI acts as a senior developer guide to codebases it can read.}

\qa{Can AI help with open source contributions?}
{Yes.
Paste the issue description and the relevant code, and ask: ``How would I implement
a fix for this issue?''
Ask AI to review whether your PR matches the project's coding style.
AI can also explain unfamiliar code in open-source projects quickly.}

\qa{How do I use AI for dependency management?}
{Ask: ``What are the risks of upgrading [package] from version X to version Y?
What breaking changes should I check for?''
Cross-reference with the package's official changelog --- AI may not have the latest
release notes.
Ask: ``Are there better-maintained alternatives to this dependency?''}

\qa{Can AI write commit messages?}
{Yes --- paste your diff and ask for a commit message in Conventional Commits format.
AI extracts the ``what'' of the change reliably.
You add the ``why'' if the business context is not obvious from the code.
Review for accuracy: AI sometimes describes what the code does rather than why it
was changed.}

\qa{How do I use AI to write release notes?}
{Paste the git log or PR titles from the release.
Ask: ``Summarise these changes as user-facing release notes. Group by feature, fix,
and breaking change. Avoid technical jargon.''
Edit for your audience and verify that nothing customer-impacting was omitted.}

\qa{How do I avoid over-reliance on AI?}
{Notice when you feel unable to start a task without AI assistance first.
Run periodic sessions where you deliberately solve problems without AI.
Ensure you understand every piece of AI-generated code in your codebase well enough
to debug it when it breaks in production.}

\qa{Can AI help with compliance requirements?}
{AI can explain GDPR, HIPAA, SOC 2, and PCI-DSS concepts and suggest implementation
patterns.
It can generate compliance documentation drafts and control descriptions.
It cannot give legal advice.
All compliance decisions must be reviewed and signed off by a qualified compliance
professional.}

\qa{How do I use AI for algorithm selection?}
{Describe the problem: data size, input characteristics, performance requirements,
and whether interpretability matters.
Ask for three algorithmic approaches with complexity analysis and tradeoffs.
Ask: ``When does this algorithm degrade and what is the worst case?''}

\qa{Can AI help me estimate story points?}
{Paste the ticket description and ask: ``What tasks are involved, what are the
technical risks, and what dependencies are implied?''
Use AI's breakdown to inform your estimate --- it does not know your team's velocity,
technical debt context, or current sprint load, so you adjust accordingly.}

\qa{How do I use AI for interview preparation?}
{Ask: ``Give me 10 hard interview questions on [topic] with detailed answers.''
Ask: ``What follow-up questions would an interviewer ask for this answer?''
Simulate a coding interview by asking AI to pose a system design question and then
critique your solution.}

\qa{Can AI help with code migration?}
{AI excels at format migrations (Python 2 to 3, jQuery to vanilla JS, REST to GraphQL).
Ask for a list of things to verify manually alongside the translated code.
Test thoroughly --- subtle behavioural differences between equivalent code in two
paradigms are common.}

\qa{How do I use AI for error handling patterns?}
{Describe your scenario and ask: ``What is the recommended error handling strategy
here?''
Ask for patterns like retry with exponential backoff, circuit breaker, and dead
letter queues by name if they apply.
Ask: ``What failure modes am I not handling in this code?''}

\qa{Can AI help me understand third-party SDK code?}
{Paste the relevant SDK source or documentation section and ask for a plain-English
explanation with a usage example.
Ask: ``What is the caller responsible for and what does the SDK handle?''
Faster than reading full SDK documentation, but always test the example.}

\qa{How do I use AI for API design?}
{Describe your resource and operations.
Ask: ``Design a RESTful API for this resource following best practices. Explain naming
conventions, status code choices, and pagination strategy.''
Ask: ``What versioning strategy is appropriate for this API and why?''}

\qa{Can AI help with monitoring and observability setup?}
{Ask: ``What metrics should I instrument for a [service type] running in Kubernetes?''
AI can generate Prometheus metrics code, Grafana dashboard JSON, and alert rules
from verbal descriptions.
Threshold values and SLO definitions are yours to determine based on your system's
actual behaviour.}

\qa{How do I use AI to write better variable and function names?}
{Paste the code and ask: ``Suggest better names for these variables and functions.
Explain your reasoning for each.''
Good naming comes from AI pattern recognition of common idioms.
Reject any name that doesn't clearly communicate intent to a reader who has never
seen the code.}

\qa{Can AI help with web accessibility (a11y) in code?}
{Ask: ``Review this HTML/JSX for WCAG 2.1 AA accessibility issues.''
AI reliably identifies missing alt text, incorrect ARIA roles, keyboard navigation
gaps, and colour contrast issues.
It cannot run automated accessibility audits --- use axe or Lighthouse for that.}

\qa{How do I use AI when I'm completely stuck on a problem?}
{Describe the problem to AI as if explaining it to a colleague who knows nothing
about your system.
Ask: ``What are five possible causes of this behaviour?''
The act of constructing a complete explanation often reveals the answer.
AI's follow-up questions expose assumptions you were not examining.}

\qa{Can AI help write load testing scripts?}
{Yes.
Describe your endpoint, expected load, and any authentication flow.
Ask for a k6, Locust, or JMeter script and explain the ramp-up strategy.
Review the script for realistic user behaviour --- AI sometimes generates unrealistic
constant-load patterns that don't represent production traffic.}

\qa{How do I translate code from one language to another using AI?}
{Paste the source code and say: ``Translate this [language A] to [language B].
Maintain the same logic. Note any semantic differences.''
Ask for a list of things to manually verify after translation.
Test the translated version against the same test suite as the original.}

\qa{How do I use AI to generate documentation?}
{Paste a function and ask for a docstring with parameters, return value, exceptions,
and one example.
Ask for README sections by name (installation, configuration, usage, examples).
Edit for tone and accuracy --- AI documentation can be generic and miss the
``why'' behind design decisions.}

\qa{What should I absolutely never use AI for?}
{Never use AI to generate credentials, secrets, or cryptographic keys.
Never commit AI output without understanding it.
Never use AI as the final word on security, legal, or financial decisions.
Never share PII, customer data, or trade secrets with non-enterprise AI tools.}

\qa{How do I use AI for feature flag design?}
{Ask: ``What is the cleanest way to implement feature flags in [framework] with these
requirements: [list]?''
AI knows patterns from LaunchDarkly, Unleash, and custom implementations well.
Ask for a comparison of client-side vs server-side evaluation and the tradeoffs.}

\qa{Can AI help me understand compiler or runtime errors?}
{Always paste the full error message and complete stack trace --- not just the last line.
Ask: ``Explain this error, where it originates, and what the most common causes are.''
AI is excellent at translating cryptic compiler messages into plain language and
pointing to the root location in the call stack.}

\qa{How do I use AI for design patterns?}
{Describe your problem, not the pattern you think you need.
\textit{``I need to add behaviour to objects at runtime without modifying their class.''}
Let AI recommend the pattern with a justification.
Ask for a concrete example in your language alongside a non-example so you understand
when not to use it.}

\qa{How do I use AI for prototype development?}
{State upfront: \textit{``This is a prototype. Prioritise working code over robustness,
error handling, and tests.''}
AI is excellent for building fast functional prototypes.
Expect to rewrite significant portions before it reaches production quality ---
and plan for it from the start.}

\qa{Can AI help me design a CI/CD pipeline?}
{Describe your workflow: trigger, test, build, scan, deploy, rollback conditions.
Ask for a GitHub Actions or GitLab CI YAML and explain each stage.
Ask: ``What are the failure modes of each step and how should I handle them?''
Review generated YAML carefully --- secrets handling and environment promotion logic
are commonly misgenerated.}

\qa{How do I use AI for code style and formatting advice?}
{Paste the code and ask: ``Does this follow [PEP 8 / Google Java Style / standard Go
formatting]? List violations and how to fix them.''
For automation, use a linter (ruff, eslint, gofmt) rather than relying on AI --- linters
are deterministic and don't hallucinate rule violations.}

\qa{Can AI help with microservice decomposition?}
{Describe your domain and the monolith's main responsibilities.
Ask: ``How would you decompose this into microservices using DDD bounded contexts?
What are the coupling risks?''
AI generates a useful starting decomposition; your architects validate against actual
deployment constraints, team structure, and data ownership requirements.}

\qa{Can AI help me write better error messages for users?}
{Yes.
Paste your current error messages and ask: ``Rewrite these to be more helpful.
Each message should tell the user what went wrong, why, and what they can do next.''
Good error messages reduce support tickets significantly and AI is very good at
generating them.}

\qa{What is the future of AI in software development?}
{AI agents will increasingly write, run, and test code autonomously, with humans
as reviewers and goal-setters rather than primary authors.
The skill shifts from writing code to knowing what good code looks like and directing
AI toward it.
The developers who thrive will be those who learn to evaluate, test, and guide AI
rather than compete with it.
Start building that skill now.}
